\documentclass[12pt,a4paper]{article}
\usepackage[utf8]{inputenc}
\usepackage[T2A]{fontenc}
\usepackage[russian]{babel}
\usepackage{amsmath,amssymb}
\usepackage{geometry}
\usepackage{hyperref}
\geometry{margin=2.5cm}

\title{Единая модель электрона как прецессирующего магнитного\\
диполя: мыльная плёнка границы действия, кулон, гравитация\\
и спин-обменное взаимодействие\\
\large Подробное изложение для читателя, впервые знакомящегося с темой}
\author{А.~Дроздов; расчёты и изложение совместно с (Claude Sonnet 5)}
\date{\today}

\begin{document}
\maketitle
\tableofcontents

\begin{abstract}
Работа излагает альтернативную теоретическую модель, в которой
электрический заряд не является фундаментальной сущностью, а
возникает как кинематический эффект от прецессии пары связанных
магнитных монополей Дирака. Ключевой геометрический механизм модели
--- трёхмерная упругая гиперповерхность (<<мыльная плёнка>>),
образующая границу четырёхмерной области интегрирования действия
в пространстве Минковского; эта плёнка мгновенно и нелокально
подстраивает свою форму, минимизируя площадь при условии
ортогональности мировым трубкам частиц (условие Ферми). Работа
даёт полное, самодостаточное изложение теории для читателя, не
знакомого с предметом: вводится понятие мировой трубки и границы
интегрирования, выводится явный вид поправки к дифференциалу
интервала (поправка Ферми), объясняется происхождение и физический
смысл <<изотропного источника продавливания>> плёнки. Отдельно и
подробно разбираются три конкурирующих механизма возникновения
эффективного электрического кулоновского взаимодействия между двумя
электрически нейтральными (по суммарному магнитному заряду)
прецессирующими диполями, даётся количественная оценка каждого
механизма и проверяется, способен ли какой-либо из них одновременно
объяснить и кулон, и ньютоновскую гравитацию. Отдельный раздел
посвящён гипотезе о механизме спин-обменного взаимодействия,
вытекающей из той же геометрической картины.
\end{abstract}

\section{Введение: зачем нужна эта модель}

Стандартная физика описывает электрон как точечную частицу с
электрическим зарядом $e$ и собственным моментом (спином), природа
которого остаётся во многом формальной: спин вводится как внутреннее
квантовое число, а не как результат какого-то реального вращения или
движения. Попытка представить спин электрона как результат обычного
механического вращения заряженного шарика классического радиуса
наталкивается на известную трудность: чтобы получить наблюдаемый
магнитный момент (магнетон Бора), поверхность такого шарика должна
была бы вращаться со скоростью, в несколько сотен раз превышающей
скорость света --- то есть классическая картина спина как вращения
электрического заряда физически невозможна.

Модель, излагаемая в этой работе, предлагает разрешение этой
трудности за счёт смены самой онтологии заряда: электрический заряд
объявляется не первичной, а <<кажущейся>> (эффективной) величиной,
возникающей из динамики двух связанных \textbf{магнитных} монополей
Дирака, которые прецессируют друг вокруг друга. Поскольку вращается
не электрический, а магнитный заряд, ограничение <<скорость меньше
$c$>> оказывается совместимым с наблюдаемой величиной магнитного
момента.

Второй ключевой элемент модели --- геометрическая интерпретация всех
видов взаимодействия (кулоновского, гравитационного, спин-обменного)
не через силовые поля в духе Максвелла или искривление
пространства-времени в духе Общей теории относительности (ОТО), а
через динамику упругой гиперповерхности --- <<мыльной плёнки>>,
образующей границу интегрирования действия. Это направление восходит
к малоизвестной работе Энрико Ферми 1923 года о парадоксе
электромагнитной массы \cite{Fermi1923}, где рассматривалось смещение
границы интегрирования действия вблизи ускоряющегося заряда, --- и к
калибровочной теории дислокаций в кристаллах, где магнитно-подобные
заряды возникают как топологические дефекты решётки
\cite{Lazar2010,KadicEdelen1983}.

Работа организована так, чтобы читатель, не знакомый ни с одной из
этих трёх составляющих (спин электрона, теория дислокаций,
геометрия границы действия), мог последовательно понять всю
конструкцию.

\section{Мировые трубки, граница интегрирования действия и мыльная плёнка}

\subsection{Что такое мировая трубка}

В теории относительности история движения точечной частицы во
времени изображается \emph{мировой линией} --- кривой в
четырёхмерном пространстве-времени (три пространственные координаты
плюс время). Если частица имеет конечный (не нулевой) размер, то её
история во времени --- это не линия, а трёхмерная <<труба>>,
протянутая вдоль направления движения: в каждый момент собственного
времени поперечное сечение этой трубы --- это сам объём частицы.
Такую фигуру называют \emph{мировой трубкой}.

В рассматриваемой модели у каждого электрона --- не одна, а
\textbf{две} мировые трубки: одна для монополя $+q_m$, другая для
монополя $-q_m$, закрученные друг вокруг друга по спирали вследствие
прецессии. Соответственно, для системы из двух электронов (например,
двух атомов позитрония) в четырёхмерном пространстве-времени
присутствуют четыре мировые трубки одновременно.

\subsection{Действие и его область интегрирования}

В классической и релятивистской механике поведение системы находится
из принципа наименьшего действия: среди всех мыслимых траекторий
движения реализуется та, для которой величина, называемая действием
$S$, минимальна (точнее, стационарна). Действие вычисляется как
интеграл некоторой плотности (лагранжиана) по четырёхмерному объёму
пространства-времени:
\begin{equation}
S = \int_\Omega \mathcal L\, d^4x.
\end{equation}
В обычной теории поля область интегрирования $\Omega$ --- это либо
всё пространство-время, либо некий фиксированный, заранее заданный
объём, и об этом обычно не задумываются специально.

Принципиальная новизна рассматриваемой модели состоит в том, что
область интегрирования $\Omega$ --- не фиксированная, а
\textbf{динамическая} величина, определяемая отдельной геометрической
задачей. Граница этой области, трёхмерная гиперповерхность
$\Sigma=\partial\Omega$, и есть та самая <<мыльная плёнка>>. Название
выбрано по прямой физической аналогии: как обычная мыльная плёнка,
натянутая на проволочный каркас, принимает форму минимальной площади
между границами каркаса, так и эта гиперповерхность принимает форму
минимальной трёхмерной <<площади>> (3-объёма) между мировыми
трубками, которые она обхватывает, как каркас.

\subsection{Условие Ферми: ортогональность плёнки мировым трубкам}

Форма плёнки не произвольна: она подчиняется дополнительному
геометрическому требованию, восходящему к Ферми \cite{Fermi1923} ---
в каждой точке, где плёнка пересекает мировую трубку частицы, плёнка
должна быть \textbf{строго ортогональна} 4-скорости этой частицы в
данной точке. Это требование называется \emph{условием Ферми} и в
формулировке автора модели описывается как то, что плёнка <<скользит
вдоль оси направления интервалов $ds$ каждой мировой трубки, сохраняя
как нормальность трубкам, так и минимизируя свою 3-поверхностную
энергию>> \cite{coax}.

Итоговая геометрическая задача, таким образом, состоит из двух
одновременных требований:
\begin{enumerate}
\item плёнка минимизирует свою собственную (трёхмерную) площадь;
\item в точках пересечения с каждой из мировых трубок плёнка строго
ортогональна 4-скорости этой трубки в данной точке.
\end{enumerate}
Это в точности вариационная задача типа минимальной поверхности
(математический аналог мыльной плёнки, натянутой на проволочный
каркас) с дополнительными краевыми условиями на ориентацию поверхности
в точках контакта с каркасом.

\subsection{Нелокальный, мгновенный характер подстройки плёнки}

Принципиально важное свойство модели, которое нужно подчеркнуть
отдельно, поскольку оно резко отличает её от привычной физики полей:
подстройка формы плёнки происходит \textbf{нелокально и мгновенно} по
всему пространству одновременно, а не распространяется с конечной
скоростью (как, например, электромагнитное поле или гравитационные
волны). В формулировке автора модели: плёнка мгновенно минимизирует
свою поверхностную энергию, и при бесконечно малом изменении фазы
процесса скользит вдоль направления интервалов $ds$ каждой мировой
трубки \cite{coax}. Иначе говоря, если где-то в пространстве меняется
ускорение одной из частиц, форма плёнки перестраивается по всему
объёму сразу же, без запаздывания на распространение сигнала ---
ровно так, как классическая (идеальная, безынерционная) мыльная
плёнка в статике мгновенно принимает новую равновесную форму при
смещении проволочного каркаса.

\subsection{Ключевое отличие от ОТО и от <<браных>> теорий: как упругая энергия плёнки входит в задачу}

Здесь необходимо явно устранить возможное недопонимание, часто
возникающее по аналогии с другими известными теориями.

В Общей теории относительности гравитация описывается через
искривление самого пространства-времени: метрический тензор
искривляется под действием тензора энергии-импульса материи, и это
искривление входит непосредственно в действие Гильберта--Эйнштейна
как отдельный (кривизнный) член, складываемый с материальным
действием. В теориях типа мембран/бран в физике высоких энергий
поверхностное натяжение границы также \textbf{прибавляется} к
основному действию как дополнительное слагаемое, пропорциональное
площади границы.

В рассматриваемой модели --- принципиально другая логика. Автор
модели формулирует это как поправку к предыдущей, ошибочной попытке
смешать эти два подхода: упругая энергия плёнки --- колоссальная
величина, но она \emph{не} является частью общего действия; вместо
этого воздействие плёнки на вещество происходит через хронометрический
механизм, тогда как сама природа интегрируемой силы берётся из
обычного 4-объёмного лагранжиана \cite{coax}.

Иначе говоря: \textbf{упругая энергия плёнки НЕ прибавляется отдельным
слагаемым к интегралу действия}. Вместо этого форма плёнки (и,
следовательно, минимизация её упругой энергии) находится как
\emph{отдельная, предварительная} геометрическая задача --- задача о
минимальной поверхности с условием ортогональности Ферми, --- и эта
задача решается \emph{до} вычисления самого физического (материального)
действия, а не одновременно с ним. Результат этой предварительной
геометрической задачи используется затем только одним, узким
способом: он определяет \textbf{локальный вес} (масштабный множитель)
для дифференциала интегрирования исходного, обычного, объёмного
лагранжиана $\mathcal L$ по каждой мировой трубке. Сама упругая
природа плёнки, таким образом, скрыта не в дополнительном слагаемом
действия, а в геометрической форме области интегрирования, которая
предопределяет, каким весом берётся вклад каждой точки трубки в
итоговый физический интеграл. Численно эта упругость характеризуется
астрономически большой величиной жёсткости вакуума $\gamma_{3D}$
(планковское давление, см. раздел~\ref{sec:model}), но входит в
конечный результат только через геометрию, а не как отдельный
энергетический член.

\subsection{Поправка к дифференциалу интервала: формула Ферми}
\label{sec:fermi-correction}

Теперь выведем явный вид той самой поправки к дифференциалу
интегрирования, о которой шла речь выше. Возьмём одну мировую трубку
(один заряд), движущуюся с ускорением. Пусть $O$ --- произвольная
точка на этой трубке (например, момент, когда скорость заряда равна
нулю), а $dt_0$ --- элементарный интервал времени в этой точке. Нас
интересует, каким будет соответствующий элементарный интервал $dt$ в
некоторой другой точке $P$, смещённой от $O$ на вектор
$(P-O)$, --- то есть насколько <<толщина>> слоя плёнки (расстояние
между соседними срезами) отличается от точки к точке.

Точный результат Ферми (для общего случая, когда скорость в точке $O$
не равна нулю) записывается через вектор кривизны мировой линии
$\vec K$:
\begin{equation}
dt = dt_0\Big[1-\vec K\cdot\overrightarrow{(P-O)}\Big],
\end{equation}
где $\vec K$ направлен по главной нормали к мировой линии и имеет
модуль, равный её кривизне \cite{Fermi1923}. Для важного частного
случая, когда скорость в точке $O$ равна нулю (то есть $O$ --- момент
поворота или максимального отклонения в колебательном/вращательном
движении), кривизна мировой линии выражается непосредственно через
собственное (3-мерное) ускорение $\vec\Gamma$ как $\vec
K=-\vec\Gamma/c^2$, и формула принимает более простой и наглядный вид
\cite{Fermi1923}:
\begin{equation}
\boxed{\ dt = dt_0\left(1+\frac{\vec\Gamma\cdot\overrightarrow{(P-O)}}{c^2}\right)\ }
\end{equation}
Физический смысл этой формулы прост: если точка $P$ находится
\textbf{впереди} по направлению ускорения (то есть
$\vec\Gamma\cdot\overrightarrow{(P-O)}>0$), локальный темп времени в
этой точке (а значит, и толщина соответствующего слоя плёнки)
\textbf{увеличивается} относительно $O$; если $P$ находится позади по
направлению ускорения --- \textbf{уменьшается}. Это и есть буквальный
геометрический смысл выражения <<плёнка продавливается>>: не
физическое проминание какого-то материального листа, а локальное
изменение темпа <<нарезки>> границы интегрирования по времени вблизи
ускоряющегося заряда.

Ферми показывает, что вклад такого (переменного по толщине) слоя в
интеграл действия для точечного заряда компенсирует расходящийся член
и превращает исходный (наивно вдвое завышенный из-за <<парадокса
4/3>>) коэффициент электромагнитной массы в правильный
\cite{Fermi1923}. Именно эта конструкция --- локальная поправка к
дифференциалу интегрирования, получаемая из геометрии границы, а не
из отдельного силового поля, --- берётся в рассматриваемой модели как
образец и обобщается на систему многих взаимодействующих (магнитных)
зарядов.

\subsection{От одной трубки к системе трубок: изотропный источник продавливания}
\label{sec:isotropic-source}

Формула Ферми, приведённая выше, --- \textbf{локальное} утверждение
для одной трубки: она говорит, как меняется толщина слоя плёнки в
малой окрестности одной ускоряющейся частицы. Чтобы описать
взаимодействие \textbf{нескольких} трубок сразу (например, для
вычисления гравитационного или кулоновского притяжения между двумя
атомами), нужно перейти от локальной формулы к \textbf{глобальному
полю} --- скалярной функции <<веса>> $w(x)$, заданной во всём
пространстве и удовлетворяющей во всём объёме уравнению типа
Лапласа--Пуассона, аналогичному уравнению электростатики, но не для
электрического потенциала, а для веса дифференциала интегрирования.
На каждой отдельной мировой трубке это глобальное поле обязано в
точности воспроизводить локальную формулу Ферми --- это и есть
граничное условие такой задачи.

Здесь возникает ключевой практический вопрос: чему равен
\textbf{источник} этого уравнения --- то есть какая физическая
величина, характеризующая ускоряющуюся частицу, <<питает>> поле веса
$w(x)$, подобно тому как электрический заряд питает электростатический
потенциал?

Прямая (наивная) идея --- взять сам вектор собственного ускорения
$\vec\Gamma$, как он стоит в формуле Ферми выше. Однако у такого
источника есть решающий недостаток: он \textbf{направленный} (векторный),
а значит, при суммировании по всем направлениям вокруг частицы (то есть
при усреднении по замкнутой поверхности, окружающей частицу, или по
периоду вращения/прецессии) он в точности гасится --- по той же
причине, по которой сумма всех единичных векторов нормали к сфере равна
нулю. Формально: поток направленной величины $\vec\Gamma\cdot\hat n$ по
замкнутой сфере всегда равен нулю независимо от направления
$\vec\Gamma$. Значит, такой источник в принципе не может дать
результирующее, не усредняющееся в ноль притяжение или отталкивание
между удалёнными объектами --- он способен объяснить лишь быстро
затухающий (дипольного типа) эффект.

Чтобы получить эффект, который \textbf{не} гасится при усреднении по
направлению или по времени (то есть настоящее, устойчивое
взаимодействие между удалёнными частицами), источник должен быть
\textbf{изотропным} --- то есть не зависящим от направления ускорения,
а зависящим только от его величины. Простейшая такая величина ---
\textbf{квадрат} вектора ускорения. Отсюда и вводится <<изотропный
источник продавливания>>:
\begin{equation}
\boxed{\ \sigma_i = \frac{m_i\,(\vec W_i\cdot\vec W_i)}{c^2}\ }
\end{equation}
где $\vec W_i$ --- собственное (в системе отсчёта, связанной с
самой частицей) ускорение $i$-й мировой трубки, а $m_i$ --- её масса.
Физический смысл этой величины: это мера того, насколько сильно и
<<равномерно во все стороны>> данная частица деформирует (продавливает)
границу интегрирования вокруг себя, --- в отличие от направленной
формулы Ферми, которая описывает лишь \emph{наклон} границы в данном
направлении. Поскольку $\sigma_i\ge0$ всегда (это квадрат, а не сам
вектор), такой источник не зависит от знака или направления ускорения
и поэтому не гасится ни при усреднении по прецессии частицы, ни при
усреднении по направлению наблюдения.

Важное следствие для модели прецессирующего магнитного диполя,
рассматриваемого как модель электрона: даже когда центр масс электрона
целиком покоится (движение отсутствует), каждый из двух монополей
внутри него продолжает двигаться по кругу вследствие прецессии, а
значит обладает ненулевым центростремительным ускорением
$W_i\sim\omega^2a$ (где $\omega$ --- частота прецессии, $a$ --- её
радиус). Отсюда следует, что источник $\sigma_i$ никогда не обращается
в ноль, даже для покоящейся частицы, --- то есть источник
взаимодействия связан не с движением центра масс, а с внутренней,
неустранимой динамикой структуры частицы. Это свойство прямо
согласуется с принципом эквивалентности: реальная гравитация тоже не
исчезает, когда тело покоится или движется равномерно.

Решение уравнения Лапласа--Пуассона с таким (изотропным, точечным)
источником даёт стандартный вид потенциала, убывающего как $1/R$:
\begin{equation}
\delta w(R) = -\frac{\sigma_i}{\gamma_{3D}\,R},
\end{equation}
где $\gamma_{3D}$ --- жёсткость (модуль упругости) вакуумной плёнки, о
которой пойдёт речь в разделе~\ref{sec:model}. Именно эта конструкция
и лежит в основе Варианта~2 расчёта, подробно разбираемого в
разделе~\ref{sec:variant2}.

\section{Калибровочная теория дислокаций Лазара (T(3)-теория)}
\label{sec:lazar}

Прежде чем перейти к трём вариантам расчёта кулона, необходимо
изложить теоретическую базу, на которую опирается наиболее сложный
(нелинейный) из трёх вариантов --- калибровочную теорию дислокаций,
разработанную М.~Лазаром \cite{Lazar2010}.

\subsection{Общая идея}

Континуальная теория дислокаций описывает кристалл со множеством
линейных дефектов (дислокаций) как упругую среду, обладающую
дополнительными степенями свободы, связанными с несовместностью
(некомпатибельностью) деформации. Лазар строит эту теорию как
калибровочную теорию группы трёхмерных трансляций $T(3)$, по прямой
аналогии с калибровочной теорией электромагнитного поля Максвелла.

Основные состояния системы --- физическая скорость $v^a$ и упругая
дисторсия (искажение) $\vartheta^a$ материального континуума. Локальная
трансляционная симметрия требует введения двух калибровочных
потенциалов: векторной 1-формы $\varphi^a$ и векторной 0-формы $\phi^a$,
которые компенсируют нарушение инвариантности при локальных сдвигах.
Гейдж-инвариантные (наблюдаемые) величины строятся как
\begin{equation}
\vartheta^a = d\xi^a+\varphi^a,\qquad v^a=\dot\xi^a+\phi^a,
\end{equation}
где $\xi^a$ --- поле деформации (аналог поля Голдстоуна).

\subsection{Полевые напряжённости: плотность и ток дислокаций}

Из калибровочных потенциалов строятся две базовые напряжённости --
прямые аналоги магнитного и электрического полей Максвелла:
\begin{align}
T^a &= d\varphi^a \quad(\text{плотность дислокаций, аналог } \vec B),\\
I^a &= d\phi^a-\dot\varphi^a \quad(\text{дислокационный ток, аналог } \vec E).
\end{align}
Величина $T^a$ характеризует, насколько плотно упакованы линейные
дефекты (дислокации) в данной точке и в каком направлении они
ориентированы, --- физически это прямой аналог того, что в
электромагнетизме называется магнитным полем: интегрирование $T^a$ по
замкнутому контуру, обхватывающему дефект, даёт квантованный вектор
Бюргерса --- топологическую характеристику дефекта, полностью
аналогичную квантованному магнитному заряду монополя Дирака
\cite{Lazar2010}. Именно это соответствие и служит математическим
мостом между теорией дислокаций в кристаллах и моделью электрона как
пары магнитных монополей.

Величина $I^a$ характеризует \emph{ток} дефектов --- скорость их
движения через данную точку пространства, и по построению
пропорциональна градиенту физической скорости среды \cite{Lazar2010};
это прямой аналог электрического поля в теории Максвелла.

Обе величины удовлетворяют тождествам, аналогичным закону Гаусса для
магнитного поля и закону электромагнитной индукции Фарадея.

\subsection{Возбуждения и определяющие соотношения}

Для замыкания теории вводятся отклики (excitations), канонически
сопряжённые напряжённостям: псевдомоментное напряжение $H_a$ и поток
дислокационного импульса $D_a$ --- прямые аналоги $\vec H$ и $\vec D$:
\begin{equation}
H_a := -\frac{\partial L_{disl}}{\partial T^a},\qquad
D_a := \frac{\partial L_{disl}}{\partial I^a}.
\end{equation}
Дополнительно вводятся дислокационное напряжение $E_a$ и
дислокационный импульс $\pi_a$ --- величины, которые, как станет ясно
ниже, отвечают именно за нелинейность теории:
\begin{align}
E_a &:= e_a\lrcorner L_{disl} + (e_a\lrcorner T^b)\wedge H_b
- (e_a\lrcorner I^b) D_b,\\
\pi_a &:= (e_a\lrcorner T^b)\wedge D_b = -(e_a\lrcorner D_b)\wedge T^b.
\end{align}

\subsection{Полевые уравнения типа Янга--Миллса и физический смысл нелинейности}
\label{sec:nonlinearity-explained}

Вариация полного лагранжиана по гейдж-инвариантным переменным даёт
замкнутую систему из двенадцати уравнений \cite{Lazar2010}:
\begin{align}
dD_a - \pi_a &= p_a, \label{eq:28}\\
dH_a - \dot D_a - E_a &= \Sigma_a. \label{eq:29}
\end{align}
где $p_a$ --- упругий импульс, а $\Sigma_a$ --- упругое напряжение
(внешние, <<упругие>> источники, не связанные с самими дефектами).

Здесь необходимо подробно, без жаргона, объяснить, что именно означает
слово <<нелинейность>> в данном контексте, поскольку именно оно лежит
в основе наиболее интересного (Варианта~3) механизма кулона.

В классической электродинамике Максвелла уравнения линейны: поле,
создаваемое одним зарядом, никак не зависит от присутствия других
зарядов и других полей --- поля просто складываются (принцип
суперпозиции). Именно поэтому, например, два световых луча,
пересекаясь в пустоте, не рассеиваются друг на друге. Уравнения
электромагнетизма Максвелла --- \emph{абелева} калибровочная теория:
объект, переносящий взаимодействие (фотон), электрически нейтрален и
не взаимодействует сам с собой.

В неабелевых калибровочных теориях (классический пример --- теория
сильного взаимодействия, квантовая хромодинамика) ситуация другая:
переносчики взаимодействия (глюоны) сами обладают <<цветовым
зарядом>> --- тем же типом заряда, который они переносят между
кварками, --- и поэтому взаимодействуют друг с другом. Это приводит к
тому, что полевые уравнения становятся \textbf{нелинейными}: поле не
просто складывается из вкладов отдельных источников --- сами поля
входят в правую часть уравнений как дополнительные источники для
самих себя.

Именно это и происходит в уравнениях (\ref{eq:28})--(\ref{eq:29})
теории Лазара: величины $\pi_a$ и $E_a$, определённые выше,
\textbf{сами выражаются} через основные поля $T^a$ и $I^a$ (через
отклики $H_a,D_a$), и при этом выступают в правой части уравнений как
дополнительные источники наряду с внешними $p_a,\Sigma_a$. Иными
словами, калибровочные поля дислокаций порождают собственные
источники, тем самым внося вклад в свои же собственные упругие
источники \cite{Lazar2010} --- ровно та ситуация самодействия, которая
характерна для неабелевых теорий типа Янга--Миллса, а не для линейной
электродинамики. Именно поэтому уравнения (\ref{eq:28})--(\ref{eq:29})
явно называются <<нелинейными>> и <<уравнениями типа
Янга--Миллса>> \cite{Lazar2010}.

Из этой системы выводится динамическая сила Пича--Кёлера --- прямой
аналог силы Лоренца:
\begin{equation}
f_a = -\dot\pi_a - dE_a = \dot p_a + d\Sigma_a
= (e_a\lrcorner I^b)\,p_b + (e_a\lrcorner T^b)\wedge \Sigma_b.
\end{equation}

\subsection{Линейная (квадратичная) версия теории}

Для локального, линейного, изотропного континуума определяющие
соотношения принимают вид
\begin{equation}
H_a = a_1\,g_{ab}\,{}^\star T^b,\qquad D_a = f_1\,g_{ab}\,{}^\star I^b,
\end{equation}
с константами $a_1$ (размерность силы) и $f_1$ (размерность
масса/длина). Лагранжиан дислокаций в этом приближении билинеен:
\begin{equation}
L_{disl}=\tfrac12 I^a\wedge D_a - \tfrac12 T^a\wedge H_a.
\end{equation}
Именно эта \textbf{линеаризованная} версия теории --- в которой поля
$\pi_a,E_a$ (нелинейные, самодействующие члены) для простоты
\textbf{опускаются} --- используется в работе Брагинского для
получения статической силы Пича--Кёлера между дислокациями
\cite{Braginsky2018,Braginsky2016}.

\section{Применение Брагинского: вывод сил из принципа наименьшего действия}
\label{sec:braginsky}

Брагинский показал, что силы, полученные Кадичем и Эделеном в
континуальной теории дислокаций по аналогии с электродинамикой, можно
вывести непосредственно из принципа наименьшего действия для тензорного
поля взаимодействия --- 4-тензора дисторсии $\lambda_i,A_{ij}$
\cite{Braginsky2018}. В его формулировке из минимума действия
получается сила, представляющая собой сумму двух слагаемых:
\begin{equation}
f_{ij}=p_i\lambda_j - e_{ijmn}p_mv_n,
\end{equation}
где $p_i$ --- импульс, $v_n$ --- скорость потока \cite{Braginsky2018}.
Первое, центрально-симметричное слагаемое $p_i\lambda_j$ --- прямой
аналог кулоновской силы; второе, вихревое --- прямой аналог силы
Лоренца \cite{Braginsky2018}. Автор специально подчёркивает, что
полученные таким образом силы можно использовать в любой среде, а не
только в твёрдом теле, поскольку они выводятся из общего принципа
наименьшего действия, а не из специфики кристаллической решётки
\cite{Braginsky2018}. Именно этот вывод и служит физическим
обоснованием переноса всей конструкции с реальных дислокаций в
кристаллах на гипотетические <<дефекты>> вакуума в модели Дроздова.

Отдельная методологическая деталь, отмеченная самим Брагинским: в
отличие от теории Кадича--Эделена, где локально калибруется сама
подгруппа трансляций, в развитии Брагинского подгруппа трансляций
глобальна, а представление волновой функции --- локально
\cite{Braginsky2016}; это меняет трансформационные свойства
компенсирующего поля, но не меняет итоговую физическую структуру сил.

\section{Модель Дроздова: электрон как прецессирующий магнитный диполь}
\label{sec:model}

\subsection{Онтология вакуумной плёнки}

В модели Дроздова физический вакуум --- не пустота, а упругая среда
планковской жёсткости $\gamma_{3D}$ (планковское давление, размерность
Паскаля). Граница четырёхмерного интеграла действия образует
трёхмерную гиперповерхность (<<мыльную плёнку>>) в пространстве
Минковского, подробно описанную в разделе~2. Единственные
фундаментальные заряды --- магнитные монополи Дирака, интерпретируемые
как топологические дефекты (дислокации) решётки вакуума --- отсюда
прямая связь с калибровочной теорией дислокаций Лазара
(раздел~\ref{sec:lazar}) и её физической интерпретацией Брагинским
(раздел~\ref{sec:braginsky}).

\subsection{Структура электрона}

Электрон представляется парой связанных монополей $+q_m,-q_m$,
прецессирующих вокруг общей оси с частотой $\omega$ на радиусе $a$
(комптоновская длина волны электрона, что даёт согласие с наблюдаемым
магнетоном Бора). Электрический заряд не является фундаментальным ---
он возникает как кинематический псевдозаряд из наклона/искривления
границы интегрирования (условие Ферми, раздел~\ref{sec:fermi-correction})
вблизи движущихся монополей.

\section{Три механизма электрического кулона между двумя нейтральными диполями}
\label{sec:variants}

Постановка задачи для всех трёх вариантов одинакова: рассматриваются
два прецессирующих магнитных диполя (два <<электрона>>, например, в
составе двух атомов позитрония), расположенных на макроскопическом
расстоянии $R$ друг от друга. Каждый диполь состоит из монополей
$+q_m,-q_m$ на расстоянии $2a$ друг от друга, прецессирующих с
частотой $\omega$; разность фаз прецессии двух диполей обозначена
$\Delta\phi$. Ключевой факт, определяющий сложность задачи: суммарный
магнитный заряд каждого диполя точно равен нулю ($q_m-q_m=0$) в любой
момент времени --- то есть диполь электрически (в магнитном смысле)
полностью нейтрален. Вопрос состоит в том, может ли взаимодействие
двух таких нейтральных объектов всё же дать эффективное
\emph{электрическое} кулоновское притяжение/отталкивание, знак
которого зависит от взаимной ориентации диполей, --- то есть
воспроизвести наблюдаемый закон Кулона $U\propto \pm1/R$.

\subsection{Вариант 1: прямое (линейное) магнитостатическое взаимодействие монополей}

Это самый прямой, <<наивный>> подход: просто сложить кулоновские
(по форме) потенциалы взаимодействия между всеми четырьмя парами
монополей (по одному монополю каждого знака в каждом диполе, то есть
$2\times2=4$ пары), с правильными знаками, определяемыми произведением
зарядов каждой пары --- в точности так, как складываются
электростатические потенциалы точечных зарядов, и в точности так, как
предписывает \emph{линейная} версия теории Лазара
(раздел~\ref{sec:lazar}), на которой основан расчёт Брагинского
(раздел~\ref{sec:braginsky}).

Итоговая знаковая сумма для четырёх пар монополей раскладывается по
малому параметру $\varepsilon=a/R$ (отношение радиуса прецессии к
расстоянию между диполями) в точный ряд:
\begin{equation}
U_{signed}(R,\phi_1,\phi_2)\cdot R = q_m^2\Big[-2\varepsilon^2\cos(\phi_1-\phi_2)
-6\varepsilon^2\cos(\phi_1+\phi_2)\Big] + O(\varepsilon^4).
\end{equation}
Принципиальный результат этого расчёта: члены нулевого порядка
($\varepsilon^0$, что соответствовало бы обычному кулону $\propto1/R$)
и первого порядка ($\varepsilon^1$, что соответствовало бы дипольному
взаимодействию $\propto1/R^2$) отсутствуют в этом разложении
\textbf{тождественно, а не приближённо}. Это прямое и неизбежное
следствие теоремы Гаусса: поскольку суммарный магнитный заряд каждого
диполя равен нулю точно в каждый момент времени, полный поток
магнитного поля через любую замкнутую поверхность, охватывающую весь
диполь, тоже равен нулю точно, --- а значит, никакого монопольного
($1/R$) вклада в принципе не может возникнуть на этом уровне, ни при
какой геометрии или скорости прецессии.

Усреднение по времени (по прецессии, $\phi_1\to\phi_1+\omega t$)
оставляет единственный не исчезающий член:
\begin{equation}
\langle U_{signed}\rangle \cdot R = -2q_m^2\varepsilon^2\cos\Delta\phi
+ O(\varepsilon^4).
\end{equation}
Это \textbf{квадрупольное} взаимодействие (по аналогии с
электростатикой, где квадрупольный потенциал убывает как $1/R^3$
из-за дополнительного множителя $\varepsilon^2=a^2/R^2$ перед $1/R$).

\textbf{Вывод по Варианту 1.} Ведущий вклад --- $\propto R^{-3}$,
знаковый (зависит от $\cos\Delta\phi$), но убывающий значительно
быстрее наблюдаемого кулона. Настоящего кулона $R^{-1}$ данный
(линейный) механизм не даёт в принципе, независимо от численных
значений параметров.

\subsection{Вариант 2: изотропный источник продавливания (линейная плёнка Ферми)}
\label{sec:variant2}

Этот вариант использует не прямое магнитостатическое взаимодействие
монополей, а механизм продавливания плёнки, подробно описанный в
разделе~\ref{sec:isotropic-source}: источник $\sigma_i=m_i(\vec
W_i\cdot\vec W_i)/c^2$, где $\vec W_i$ --- собственное центростремительное
ускорение прецессии данного диполя.

\subsubsection{Происхождение и физический смысл эффекта}

Как объяснялось выше, этот источник \textbf{изотропен} (зависит от
квадрата ускорения, а не от его направления), а потому не гасится ни
при усреднении по прецессии, ни при усреднении по направлению.
Решение уравнения Лапласа--Пуассона для такого источника даёт
истинный монопольный потенциал:
\begin{equation}
\delta w(R) = -\frac{\sigma_i}{\gamma_{3D}\,R}.
\end{equation}
Принципиальная особенность этого эффекта, которую необходимо
подчеркнуть отдельно: поскольку источник $\sigma_i$ содержит
\textbf{квадрат} ускорения, он совершенно не чувствителен к знаку или
направлению вращения --- он одинаков для пары электрон--электрон,
электрон--позитрон, для любых двух прецессирующих диполей независимо
от относительной фазы их прецессии $\Delta\phi$ (на этом порядке
малости). Это принципиальное отличие от Варианта~1 и Варианта~3
(ниже): у $\sigma_i$ \textbf{нет знака}, а значит, взаимодействие,
которое он порождает, всегда одного знака (притяжение) --- в точности
так, как ведёт себя гравитация (все массы притягиваются, независимо
от их <<ориентации>>), но принципиально не так, как ведёт себя
электрический кулон (одноимённые заряды отталкиваются, разноимённые
притягиваются). Поэтому данный механизм называется
\textbf{гравитационно-подобным}, а не электрически-подобным: он
структурно повторяет свойство гравитации (всегда притяжение,
пропорциональное <<массе>>, то есть энергии внутреннего движения
частицы), но принципиально не способен воспроизвести знаковую природу
электрического взаимодействия.

\subsubsection{Численная оценка: объясняет ли Вариант 2 гравитацию}

Чтобы проверить, способен ли этот механизм количественно объяснить
\emph{саму} наблюдаемую гравитационную постоянную $G$, нужно
подставить в $\sigma_i$ реальные параметры прецессии электрона.
Радиус прецессии $a$ равен комптоновской длине волны электрона
($a=\hbar/(m_ec)\approx3{,}86\times10^{-11}$~см), а скорость прецессии
$v=\omega a$, необходимая для совпадения с наблюдаемым магнетоном
Бора, лежит в диапазоне $v\approx0{,}03c$--$0{,}09c$. Подставляя эти
числа в $\sigma_i=m_e(v^2/a)^2/c^2$ и сравнивая результат с реальным
отношением гравитационной и кулоновской силы между двумя электронами,
$Gm_e^2/e^2\approx2{,}4\times10^{-43}$, получаем расхождение порядка
$10^{40}$--$10^{50}$ раз --- то есть кинематический источник,
построенный из скорости и радиуса прецессии электрона, оказывается на
40--50 порядков величины \textbf{слабее}, чем необходимо для
воспроизведения наблюдаемой силы гравитации.

Причина этого разрыва --- не ошибка в вычислении, а фундаментальное
структурное несоответствие масштабов. Наблюдаемая слабость гравитации
количественно выражается через безразмерную величину
$\alpha_G=Gm_e^2/(\hbar c)\approx1{,}75\times10^{-45}$ --- квадрат
отношения массы электрона к планковской массе. Планковская масса
примерно в $2\times10^{22}$ раз больше массы электрона: именно это
огромное отношение и объясняет, почему гравитация настолько слабее
других сил в реальном мире. Но кинематический источник $\sigma_i$
строится только из величин, характерных для самого электрона
(скорость, комптоновская длина, масса электрона) и никак не содержит
планковского масштаба --- поэтому он не способен самопроизвольно
воспроизвести подавляющий фактор такой величины.

\textbf{Вывод по Варианту 2.} Правильная функциональная форма
(притяжение $\propto1/R$, без знака, растущее с массой) при этом
воспроизводится верно, но количественная величина эффекта --- нет:
разрыв с реальной гравитацией составляет 40--50 порядков величины.
Иными словами, Вариант~2 показывает \emph{как} мог бы работать
подобный механизм, но не даёт правильного числа для реальной
гравитационной постоянной.

\subsection{Вариант 3: нелинейная самосвязь калибровочного поля дислокаций}
\label{sec:variant3}

\subsubsection{Что такое <<нелинейная самосвязь>> в данном контексте --- развёрнутое объяснение}

Как было подробно объяснено в разделе~\ref{sec:nonlinearity-explained},
полные (не линеаризованные) уравнения Лазара (\ref{eq:28})--(\ref{eq:29})
содержат величины $\pi_a$ и $E_a$, которые сами выражаются через поля
$T^a$ (плотность дислокаций/монополей) и $I^a$ (ток дислокаций/скорость
движения монополей) и одновременно выступают дополнительными
источниками для тех же самых полей. Физически это означает следующее:
в линейной (обычной, максвелловской) картине поле, создаваемое одним
диполем, никак не <<чувствует>> присутствие поля от другого диполя ---
поля просто складываются, и, как было показано в Варианте~1, при
сложении полей от электрически нейтрального объекта (диполя с нулевым
суммарным зарядом) результирующее дальнее поле гасится с точностью до
$1/R^3$.

В \textbf{нелинейной} картине ситуация другая: поле от диполя~1 и поле
от диполя~2 не просто складываются --- они дополнительно
\textbf{перемножаются} друг с другом через квадратичные (по полям)
члены самосвязи $\pi_a,E_a$, и этот квадратичный, перемноженный вклад
не подчиняется теореме Гаусса для суммарного заряда так же просто,
как линейная сумма потенциалов. В частности, произведение двух величин
одного знака (или квадрат величины) всегда неотрицательно и не
гасится при усреднении --- в этом смысле механизм структурно похож на
изотропный источник Варианта~2 (тоже основанный на квадрате величины,
там --- ускорения), но здесь квадратичная (перемножающая) операция
применяется не к ускорению, а к самим полям $T^a,I^a$ монополей, что,
как показано ниже, сохраняет знаковую (кулоно-подобную) зависимость от
взаимной фазы прецессии.

\subsubsection{Явный вид членов самосвязи}

Подставляя линейные определяющие соотношения $H_a=a_1(\star T^a)$,
$D_a=f_1(\star I^a)$ (раздел~\ref{sec:lazar}) в определения
дислокационного напряжения и импульса, получаем явные выражения для
членов самосвязи через основные модули упругости $a_1,f_1$ --- те же
самые константы, которые определяют и линейную (Вариант~1-подобную)
часть теории, без введения каких-либо новых, отдельно подбираемых
величин:
\begin{align}
\pi_a&=(e_a\lrcorner T^b)\wedge f_1(\star I_b),\\
E_a&=e_a\lrcorner\Big[\tfrac12 f_1\, I\wedge\star I-\tfrac12 a_1\, T\wedge\star T\Big]
+a_1(e_a\lrcorner T^b)\wedge(\star T_b)-f_1(e_a\lrcorner I^b)(\star I_b).
\end{align}

\subsubsection{Применение к системе двух прецессирующих диполей}

Для системы из двух источников (полный дислокационный поток $T=T_1+T_2$
и полный ток $I=I_1+I_2$) при подстановке в квадратичные выражения выше
автоматически возникают \textbf{перекрёстные} члены вида
$T_1\wedge\star I_2+T_2\wedge\star I_1$ --- именно они, а не самодействие
каждого диполя по отдельности (которое не зависит от $R$ и представляет
собой лишь перенормировку собственной массы каждого диполя), дают
зависящий от расстояния $R$ вклад во взаимодействие между диполями.

Чтобы оценить эти перекрёстные члены количественно, нужно перенести
абстрактную геометрию линейных дислокаций на геометрию точечных
прецессирующих монополей --- это отдельный, самостоятельный шаг,
основанный на аналогии Брагинского (раздел~\ref{sec:braginsky}), а не
точное математическое тождество. При таком переносе поле $T^a$
представляется обычным дипольным полем двух точечных монополей
($+q_m,-q_m$) на расстоянии $2a$, а поле $I^a$, поскольку оно по
построению пропорционально градиенту скорости \cite{Lazar2010},
оценивается как $I^a\sim(v/c)\,T^a$, где $v=\omega a$ --- скорость
прецессии монополя.

Крест-член $T_1\wedge\star I_2$ для двух дипольных полей вычисляется
через интеграл по всему пространству от произведения потенциалов
$\Phi_{0,1}\Phi_{0,2}$, порождаемых каждым диполем (аналог энергии
взаимодействия двух перекрывающихся полей). Этот интеграл был вычислен
аналитически (через фурье-свёртку двух точечных дипольных распределений,
разделённых расстоянием $R$) и независимо проверен численно методом
Монте-Карло с точностью порядка 5\%. Для двух дипольных моментов
$\vec p_1,\vec p_2$ ($p=2q_ma$), закреплённых на плоскости, перпендикулярной
линии, соединяющей центры диполей, и усреднённых по времени (по
совместной прецессии с постоянной разностью фаз $\Delta\phi$), получена
закрытая аналитическая формула:
\begin{equation}
\boxed{\ \langle I_{TT}(R,\Delta\phi)\rangle
= \frac{\pi p^2}{R}\cos\Delta\phi,\qquad p=2q_m a\ }
\end{equation}

Это ключевой качественный результат Варианта~3: полученное
взаимодействие --- \textbf{знаковое} (пропорционально
$\cos\Delta\phi$, то есть зависит от взаимной ориентации прецессии
двух диполей ровно так, как электрический кулон зависит от знака
зарядов), \textbf{не гасится} при усреднении по прецессии (в отличие от
монопольного и дипольного членов Варианта~1) и убывает точно как
$1/R$ --- то есть впервые среди всех трёх вариантов воспроизводит
верную \textbf{функциональную форму} закона Кулона.

Полная оценка энергии нелинейного взаимодействия с учётом модуля
$f_1$ и подавляющего множителя $(v/c)$ (от связи $I^a\sim(v/c)T^a$)
записывается как
\begin{equation}
U_{nl}(R,\Delta\phi) \approx f_1\cdot\frac{v}{c}\cdot\frac{\pi p^2}{R}\cos\Delta\phi.
\end{equation}

\subsubsection{Оценка величины и открытый вопрос о модуле $f_1$}

Ключевое отличие этой оценки от прямого магнитостатического расчёта
(Вариант~1) --- отсутствие свободно подбираемой константы связи:
величина эффекта целиком определяется модулем упругости $f_1$, тем
же самым, который в общем виде уже присутствует в конститутивных
соотношениях теории Лазара (раздел~\ref{sec:lazar}). Однако конкретное
численное значение $f_1$ для физического вакуума --- в интерпретации
настоящей модели, для самой упругой плёнки границы действия, --- в
рассмотренных материалах не зафиксировано: у Лазара $f_1$ --- это
абстрактная материальная константа теории дислокаций общего вида,
применимая к любой упругой среде, а не конкретное число, привязанное
к физическому вакууму. Таким образом, свобода в этом расчёте
заключается именно в величине физического модуля $f_1$, подлежащей
отдельному определению --- в отличие от Варианта~1, где никакой
подобной подгонки в принципе быть не может (там результат равен нулю
на нужном порядке независимо от каких-либо констант), и в отличие от
Варианта~2, где сама численная величина эффекта уже полностью
определена параметрами электрона и жёсткостью $\gamma_{3D}$, но
оказывается катастрофически мала.

\subsubsection{Способен ли Вариант 3 также объяснить гравитацию}

Естественный следующий вопрос: если Вариант~3 способен дать верную
\emph{форму} кулона, можно ли подобрать параметры так, чтобы
\textbf{одновременно} объяснить и кулон, и гравитацию, --- то есть
устранить сразу оба недостатка Варианта~2?

Чтобы ответить на этот вопрос честно, нужно разделить два вклада,
которые содержатся в общей формуле для дислокационного напряжения
$E_a$ выше. Помимо перекрёстного (сигнального, $\propto\cos\Delta\phi$)
члена $\propto f_1\,T_1\wedge\star I_2$, та же формула содержит и
\textbf{самостоятельный} квадратичный член вида $a_1\,T\wedge\star T$,
взятый для полной (суммарной) плотности $T=T_1+T_2$. Этот член
включает как вклады самодействия каждого диполя по отдельности (не
зависящие от $R$, дающие лишь перенормировку собственной массы), так и
перекрёстный член $T_1\wedge\star T_2$, зависящий от $R$. В отличие от
члена $T_1\wedge\star I_2$ (который через $I\sim(v/c)T$ приобретает
знаковую $\cos\Delta\phi$-структуру, показанную выше), необходимо
отдельно проверить структуру именно члена $T_1\wedge\star T_2$: по
своей математической природе (произведение двух одинаковых по типу
полей, а не поля и тока) он структурно ближе к самодействию, то есть
к тому же типу вклада, что и изотропный источник $\sigma_i\propto W^2$
Варианта~2, --- потенциальный кандидат на роль независимого,
гравитационно-подобного (не зависящего от знака и от $\Delta\phi$)
канала внутри того же самого нелинейного механизма.

Если это так, то модель в принципе содержит \textbf{два разных
физических канала} внутри одной и той же нелинейной теории:
знаковый, кулоно-подобный канал ($\propto f_1\cdot(v/c)\cdot\cos\Delta\phi/R$,
через перекрёстный член $T$ и $I$) и безнаковый, гравитационно-подобный
канал ($\propto a_1\cdot(\ldots)/R$, через перекрёстный член $T$ и
$T$). Поскольку эти два канала управляются \textbf{разными}
физическими модулями ($f_1$ и $a_1$ соответственно), в принципе
\emph{всегда} можно подобрать пару чисел $(a_1,f_1)$ так, чтобы
одновременно точно воспроизвести и наблюдаемую величину кулона, и
наблюдаемую величину гравитации.

Здесь необходима честная методологическая оговорка. Такая
одновременная подгонка \textbf{двух независимых} параметров под
\textbf{два независимых} наблюдаемых числа --- это не физическое
предсказание, а простая калибровка (в теории с двумя свободными
параметрами всегда можно точно подобрать их так, чтобы совпасть с
двумя любыми заданными числами; это не является проверкой теории).
Содержательным, по-настоящему проверяемым результатом было бы не
совпадение (которое достигается автоматически при двух свободных
параметрах), а появление \textbf{дополнительного, независимого}
предсказания --- например, конкретного числового \emph{соотношения}
между величиной кулона и величиной гравитации, вытекающего из
структуры теории \emph{без} дополнительной подгонки $a_1$ и $f_1$
отдельно друг от друга. Получение такого соотношения (то есть
независимая, не подгоночная связь между $a_1$ и $f_1$, вытекающая,
например, из более фундаментальных свойств самой вакуумной решётки)
остаётся открытой задачей и не решено в рамках рассмотренных
материалов.

\subsection{Сводная таблица результатов}

\begin{center}
\begin{tabular}{|l|c|c|c|}
\hline
Вариант & Ведущий порядок & Знак/$\Delta\phi$? & Кулон $R^{-1}$ со знаком? \\
\hline
1. Линейная магнитостатика & $R^{-3}$ & да ($\cos\Delta\phi$) & Нет \\
2. Изотропный источник Ферми & $R^{-1}$ & нет & Нет (нет знака; не гравитация: разрыв $10^{40-50}$) \\
3. Нелинейная самосвязь & $R^{-1}$ & да ($\cos\Delta\phi$) & Да по форме; $f_1$ не откалиброван \\
\hline
\end{tabular}
\end{center}

\section{Гипотеза о механизме спин-обменного взаимодействия}
\label{sec:exchange}

\subsection{Постановка вопроса}

Помимо кулоновского и гравитационного взаимодействия, в квантовой
механике существует третий, чисто квантовый тип взаимодействия между
двумя электронами --- \emph{обменное взаимодействие}, определяющее,
например, разницу энергий между параллельной и антипараллельной
ориентацией спинов двух электронов (это лежит в основе явлений
магнетизма и химической связи). В стандартной квантовой механике это
взаимодействие возникает из требования антисимметрии волновой функции
системы одинаковых фермионов и не имеет классического
геометрического аналога.

Естественный вопрос для рассматриваемой модели: может ли геометрия
скольжения мыльной плёнки дать что-то похожее на обменное
взаимодействие --- то есть энергию, зависящую от взаимной ориентации
прецессии (а значит, в данной модели, от взаимной ориентации спинов)
двух электронов, --- без обращения к квантовой антисимметрии, а
исключительно из классической геометрии деформации общей границы
интегрирования?

\subsection{Кандидат на механизм: член более высокого порядка по $\varepsilon$}

При вычислении \textbf{безнакового} (изотропного) варианта суммы
взаимодействия всех четырёх монополей --- аналогичного тому, что
используется в оценке гравитационно-подобного канала Варианта~2, но
только теперь без учёта знаков зарядов (то есть суммы вида
$S=\sum_{i,j}1/R_{ij}$ по всем четырём парам монополей, взятой без
знаков) --- разложение по малому параметру $\varepsilon=a/R$ даёт не
только уже обсуждавшийся не зависящий от $\Delta\phi$ ведущий член, но
и следующий по порядку член четвёртой степени по $\varepsilon$, который
явно зависит от взаимной фазы прецессии по закону
$\cos(2\Delta\phi)$ (а не $\cos\Delta\phi$, как в кулоновском
Варианте~3):
\begin{equation}
S(R,\Delta\phi)\cdot R = 4+2\varepsilon^2+\varepsilon^4\left(\frac{27}{8}+\frac98\cos2\Delta\phi\right)+O(\varepsilon^6).
\end{equation}

\subsection{Физическая интерпретация}

Появление именно \textbf{удвоенного} угла $2\Delta\phi$ (а не просто
$\Delta\phi$) в этом члене физически значимо и отличает его от
кулоновского механизма. Зависимость с периодом $\pi$ по $\Delta\phi$
(а не с периодом $2\pi$) характерна для взаимодействий квадруполь--квадруполь
или, в более общем виде, для эффектов, чувствительных не к самому
направлению вращения, а к \textbf{оси} вращения без различения двух
противоположных направлений вдоль этой оси, --- то есть к тому
свойству геометрии, которое отличает <<параллельную>> ориентацию от
<<перпендикулярной>>, но не различает два противоположных направления
внутри параллельной ориентации. Это качественно (хотя не тождественно
количественно) напоминает характерное свойство настоящего квантового
обменного взаимодействия --- зависимость энергии от взаимной
ориентации спинов (в простейшей модели Гейзенберга --- от скалярного
произведения спиновых операторов $\vec S_1\cdot\vec S_2$), которая
также определяется относительной, а не абсолютной ориентацией.

Рабочая гипотеза, вытекающая из этого наблюдения, формулируется так:
если два прецессирующих диполя сближаются на расстояние, при котором
члены четвёртого и более высоких порядков по $\varepsilon$ становятся
существенными (то есть на масштабах порядка нескольких радиусов
прецессии $a$, что в реальном электроне соответствует
субкомптоновским масштабам сближения --- то есть речь идёт не о
типичном атомном/молекулярном расстоянии, где действует наблюдаемый
обмен, а о гипотетическом механизме на существенно меньших
масштабах), возникающая зависимость энергии деформации плёнки от
$\cos2\Delta\phi$ могла бы играть роль, аналогичную квантовому
обменному взаимодействию: энергия оказывается разной для случаев,
когда оси прецессии двух электронов параллельны, и когда они
перпендикулярны, --- в духе классического (не квантового по
происхождению) аналога спин-спинового взаимодействия.

\subsection{Честная оценка статуса гипотезы}

Необходимо прямо и ясно указать степень надёжности этого раздела,
чтобы не создавать ложного впечатления завершённого результата.
Приведённая выше формула для члена $\cos2\Delta\phi$ получена из
чисто геометрического (безнакового) суммирования расстояний между
монополями и \textbf{не выведена} напрямую из нелинейных уравнений
теории Лазара, в отличие от кулоновского механизма Варианта~3, где
такой прямой вывод дан явно. Кроме того, убывание этого члена по
расстоянию ($\propto\varepsilon^4/R=a^4/R^5$) --- значительно более
быстрое, чем у реального обменного взаимодействия, которое в
квантовой механике определяется перекрытием волновых функций и
убывает экспоненциально с расстоянием, а не степенным образом; прямое
количественное сравнение этих двух разных функциональных форм пока не
проведено. Таким образом, изложенное в этом разделе --- рабочая
гипотеза о возможном классическом геометрическом аналоге
спин-обменного взаимодействия, качественно мотивированная структурой
уже полученного результата, но не доведённая до количественной
проверки и не претендующая на замену квантово-механического описания
обмена.

\section{Заключение}

Изложенная модель предлагает целостную геометрическую картину, в
которой электрический заряд, гравитация и, гипотетически,
спин-обменное взаимодействие возникают из единого источника ---
динамики скольжения упругой трёхмерной границы интегрирования действия
(<<мыльной плёнки>>) вдоль мировых трубок прецессирующих магнитных
монополей. Три проверенных механизма кулоновского взаимодействия дают
разные, но взаимно дополняющие уроки: линейное магнитостатическое
взаимодействие (Вариант~1) исключает возможность настоящего кулона по
топологическим причинам (нулевой суммарный заряд диполя); изотропный
источник продавливания (Вариант~2) даёт верную форму гравитационного
взаимодействия, но количественно на 40--50 порядков слабее
наблюдаемого; и лишь нелинейная самосвязь калибровочного поля
дислокаций (Вариант~3), явно выведенная из точных уравнений Лазара без
введения произвольной константы связи, впервые воспроизводит верную
знаковую, кулоно-подобную функциональную форму взаимодействия. При
этом одновременное объяснение и кулона, и гравитации в рамках
Варианта~3 в принципе возможно за счёт двух независимых физических
модулей ($a_1$ и $f_1$), но представляет собой калибровку, а не
независимую проверку теории, до тех пор пока не будет найдено
дополнительное соотношение, связывающее эти модули без подгонки.
Гипотеза о геометрическом аналоге спин-обменного взаимодействия,
основанная на члене $\cos2\Delta\phi$ в разложении той же
геометрической конструкции, остаётся качественной и требует отдельной
количественной проверки.

\begin{thebibliography}{9}
\bibitem{Fermi1923} E. Fermi, работа 1922--1923 гг. о геометрии
интегрирования вблизи ускоряющегося заряда и парадоксе
электромагнитной массы.
\bibitem{coax} Материалы обсуждения <<Interaction of coaxial current
rings>> (рабочие заметки автора модели).
\bibitem{Lazar2010} M. Lazar, {\it On the Fundamentals of the
Three-dimensional Translation Gauge Theory of Dislocations},
Mathematics and Mechanics of Solids {\bf 16}(3), 253--264 (2010).
\bibitem{Braginsky2018} А.~Я.~Брагинский, рабочие материалы (2018).
\bibitem{Braginsky2016} А.~Я.~Брагинский, рабочие материалы (2016).
\bibitem{KadicEdelen1983} A. Kadi\'c, D.G.B. Edelen, {\it A Gauge
Theory of Dislocations and Disclinations}, Lecture Notes in Physics,
Vol. 174, Springer, Berlin (1983).
\end{thebibliography}

\end{document}
